\chapter*{引言}
\phantomsection
\addcontentsline{toc}{chapter}{引言}

量的测度这一概念是根本性的，无论在日常生活（长度、面积、体积、重量）中，还是在实验科学（电荷、磁质量等）中都是如此。如此多种多样的量的“测度”的共同特征在于：对满足某些条件的每一部分空间联结一个数，使得对两个这样的部分（假定没有公共点）的并，对应于分别赋予它们各自的数之和（测度的可加性）(*)。此外，测度通常是一个正数，这就意味着它是被测的那部分空间的递增函数 (**)。另一方面可以注意到，在实际中人们几乎从不操心去指明哪些部分空间应当看成“可测的”；当然，在任何一种测度的数学理论中，都必须毫不含糊地解决这个问题；例如，在初等几何中定义多边形的面积或多面体的体积时，人们做的就是这件事；在所有这些情形，“可测”集的族自然必须满足：其中任意两个无公共点的集合的并也是“可测的”。

在上面大多数例子中，一部分空间的测度随其直径而趋于 0：按经典的说法，一个点“没有长度”，这就是说它含于长度任意小的区间之中，因而只能赋予它长度 0；这样的量的测度称为“弥散的”。然而，力学和物理学的发展引入了这样一类量的概念：对这类量，一个尺寸可忽略的物体仍可以具有不可忽略的测度：引力的或电的“点质量”，说实在的，它们在很大程度上是数学虚构，而不是严格意义上的实验概念。于是，在数学中人们被引向考虑如下定义的测度：对一个有限集的每一点 \(a_i\) ( \(1 \leqslant i \leqslant n\) )（记此有限集为 \(F\) ）附加一个数 \(m_i\) ，即它的“质量”或“重量”，而任一集合 \(A\) 的测度，就是质量 \(m_i\) 的和，这里点 \(a_i\) 属于 \(A\) 。

与测度概念紧密联系在一起的是加权和的概念。例如，考虑空间中有限多个（引力的或电的）质量 \(m_{i}\) ，它们置于点 \(a_{i}\) （坐标为 \(x_{i}, y_{i}, z_{i}\) ）处；这组质量施加于点 b（质量为 1，坐标为 \(\alpha, \beta, \gamma\) ）上的引力沿 Oz 方向的分量（例如）是（在适当选取的单位制下的）和式

\[
\sum_ {i} m _ {i} \frac {(z _ {i} - \gamma)}{r _ {i} ^ {3}},
\]

\(r_i^2 = (x_i - \alpha)^2 + (y_i - \beta)^2 + (z_i - \gamma)^2\) 是点 \(a_i\) 与 \(b\) 之间距离的平方。换句话说，我们考虑函数

\[
f (x, y, z) = \frac {z - \gamma}{\left((x - \alpha) ^ {2} + (y - \beta) ^ {2} + (z - \gamma) ^ {2}\right) ^ {3 / 2}}
\]

在每个点 \(a_{i}\) 处的值，把它乘以该点的“权重”，再把如此得到的 f 的“加权值”全部加起来。众所周知，这样的和式在力学中不断地出现：重心和惯性矩就是最著名的例子。

如果人们想要把“加权和”的概念从点质量的情形推广到每一点测度均为零的“弥散”测度的情形，就会遇到那个外表如此悖理、并催生了积分学的问题：如何给一个具有无穷多项、而每一项单独取出来都是零的“和”赋予意义。让我们再回到计算施加于一点上的引力的例子，其中吸引的质量“连续地分布”在整个体积 V 中。如果把 V 分解为有限个（两两不相交的）子集 \(V_{i}\) ，我们就假定 V 施加于点 b 上的引力沿 Oz 方向的分量等于各个 \(V_{i}\) 施加于 b 上的引力的分量之和。而如果每个 \(V_{i}\) 的直径都很小，连续函数 \(f(x,y,z)\) 在 \(V_{i}\) 中就变化很小，于是人们被引向把 \(V_{i}\) 所施加的引力近似地看成一个点质量所会施加的引力，该点质量等于质量 \(m_{i}\) （即 \(V_{i}\) 的质量），并置于点 \(a_{i}\) 处（该点任取于体积 \(V_{i}\) 之中）。于是人们被引向取“Riemann 和” \(\sum m_{i}f(x_{i},y_{i},z_{i})\) 作为所求数的近似值；为使这在数学观点上得到论证，自然必须证明当 \(V_{i}\) 的最大直径趋于 0 时这些近似值趋于一个极限，而这由函数 f 在 V 中的一致连续性容易得到（假定 V 是紧的且点 b 不在 V 中）。

众所周知，希腊人的“穷竭法”以及用于系统地计算平面图形面积与体积的“Cavalieri 原理”，都建立在一个类似的程序上，即把所考虑的面积与体积分解成“切片”；由此得到的“加权和”不是别的，正是积分 \(\int_{a}^{b} f(x) dx\) （参见第 IV 卷第 I--III 章的历史注记）。这里同样是 f 的一致连续性蕴涵着“Riemann 和”极限的存在性；更一般地，它蕴涵类似的和式 \(\sum_{i} f(\xi_{i})(g(x_{i+1}) - g(x_{i}))\)

\((x_{i} \leqslant \xi_{i} \leqslant x_{i+1})\) （其中只假定 g 是 \([a, b]\) 上的有界递增函数）时极限存在。这一极限记作 \(\int_{a}^{b} f(x) dg(x)\) ，称为 f 关于 g 的 Stieltjes 积分，它可以看作函数 f 关于测度 \(\mu\) 的“加权和”，该测度在半开区间 \([\alpha, \beta]\) 所成的集上由公式 \(\mu([\alpha, \beta]) = g(\beta+) - g(\alpha+)\) 定义；它不再像通常的积分概念那样与微分学紧密地联系在一起 (*)。经典的“二重”积分和“三重”积分分别联系于平面面积的度量与体积的度量，情形也是如此。然而，所有这些积分概念彼此之间的联系不仅在于它们的定义，还在于下述特征：直线、平面或三维空间中某个紧子集 K 上的连续数值函数 f 的“积分” \(\mu(f)\) ，是与 K 上的连续函数空间 \(\mathcal{C}(K)\) 的元素 f 相联系的一个数；于是 \(f \mapsto \mu(f)\) 是 \(\mathcal{C}(K)\) 到 R 中的映射（有时称为“泛函”），它是：\(1^{\circ} linear\) （也就是说，\(\mu(\alpha f + \beta g) = \alpha \mu(f) + \beta \mu(g)\) 对一切标量 \(\alpha, \beta\) 和一切连续函数 f, g 成立）；\(2^{\circ} positive\) （也就是说，\(\mu(f) \geqslant 0\) 对每个满足 \(f \geqslant 0\) 的连续函数 f 成立）。

值得注意的是，反过来，这两条性质就足以刻画区间 \([a,b]\) 上的 Stieltjes 积分（F. Riesz 定理）。之所以如此，是因为从连续函数积分的值出发，能够重新构造出产生该积分的测度。这相当于（如果我们把 \(\int_{a}^{b}f(x)dx\) 解释为平面图形的面积的话），在假定积分对连续函数为已知的前提下，去计算一个区间的特征函数的积分。换句话说，问题在于以适当的方式把泛函 \(\mu(f)\) 延拓到一个既包含 \(\mathcal{C}(K)\) 又大到足以包含各区间的特征函数的函数集上。

完成这一延拓有若干种方法；其中最有意思的方法之一求助于函数空间的概念。我们知道，在空间 \(\mathbf{R}^n\) 上，范数 \(\| \mathbf{x}\|_{\infty} = \sup_{1\leqslant i\leqslant n}|x_i|\) 与 \(\| \mathbf{x}\|_1 = \sum_{i=1}^{n}|x_i|\) 定义同一个拓扑。通过“从有限过渡到无穷”，人们被引向考虑连续函数空间 \(\mathcal{C}(\mathrm{K})\) （\(\mathrm{K} = [a,b]\) 是 \(\mathbf{R}\) 中的一个紧区间）上的范数 \(\| f\|_{\infty} = \sup_{x\in \mathrm{K}}|f(x)|\) 与 \(\| f\|_1 = \int_a^b |f(x)|dx\) （在 Stieltjes 积分的情形则取 \(\int_a^b |f|dg\) ）。但在这里，这两个范数所定义的拓扑是不同的，而且对第一个范数完备的空间 \(\mathcal{C}(\mathrm{K})\) （GT, X, §1, No. 6, Th. 2 的推论 1），对第二个范数不再完备。更确切地说，可以把 \(\mathcal{C}(\mathrm{K})\) 关于范数 \(\| f\|_1\) 的完备化的元素与未必连续的函数的类等同起来，于是积分的延拓就简单地通过把泛函 \(\mu(f)\) （它定义在 \(\mathcal{C}(\mathrm{K})\) 上）连续地延拓到该空间的完备化上来实现（这一程序的技术细节在第 IV 章中陈述）。当然，我们曾假定连续函数的积分是借助一个测度（用上面回顾过的“Riemann 和”程序）来定义的；为得到 F. Riesz 定理，必须按同样的方式操作，只是把范数定义为 \(\mu(|f|)\) ，其中 \(\mu(f)\) 是定义在 \(\mathcal{C}(\mathrm{K})\) 上的正线性形式。

我们刚才概述的延拓方法不仅导出 Riesz 定理，而且它还允许对由比区间的特征函数“不连续得多”的函数所组成的类定义积分；在考察作为“可积”函数的集合的特征函数时，它同时允许把起初只对区间给出的测度延拓到相应的集合上，办法是令 \(\mu(A) = \mu(\varphi_A)\) ；这一延拓当然保持了测度的可加性与正性这两条基本性质。

前面讨论的是直线上的 Stieltjes 积分，但延拓方法可以立即搬到定义在平面或空间中的测度上，或定义在曲线上或曲面上的测度上。更一般地，在分析这些证明时可以看出，它们实际上对定义在任意局部紧空间 X 上的连续函数所成的空间 \(\mathcal{H}(X)\) 上的每一个正线性形式都有效，这些连续函数中的每一个都在某个紧子集（依具体的函数而定）之外为零。

积分理论因而可以适用的这类空间自然包括数值空间 \(R^{n}\) 以及流形；它也包括离散空间（在离散空间上，积分理论与实数的可和族的理论合而为一 (GT, IV, §7)），还包括与 R 的区间相同的、或与有限集相同的紧空间的（有限或无穷）乘积；我们以后将看到，这类乘积上的测度理论在概率演算中起着重要的作用。

把测度概念延拓到一般的局部紧空间，这在局部紧群的理论中显得特别富有成效；一般地说，每当在拓扑代数中想要“从有限过渡到无穷”，也就是说，想要把出现有限和的纯代数程序推广到“求和”必须处理无穷多项的情形，积分的概念看来就是恰当的工具。例如，我们知道 (A, III, §2)，有限群 G（在域 \(\mathbf{R}\) 上）的代数的元素，是形如 \(s \mapsto \alpha(s)\) 的映射（即 G 到 \(\mathbf{R}\) 中的映射），其乘法规则为 \(\alpha * \beta = \gamma\) ，其中 \(\gamma\) 是由下式定义的函数

\[
\gamma (s) = \sum_ {t \in G} \alpha (t) \beta (t ^ {- 1} s).
\]

对任意的局部紧群 G，这一代数的一个看来自然的推广，是 G 到 R 中的那样一些映射所成的集，它们关于 G 上某个特殊的测度 \(\mu\) （“Haar 测度”）可积，而代数中的乘法由

\[
(f * g) (s) = \int f (t) g (t ^ {- 1} s) d \mu (t).
\]

此外，一旦走上这条路，人们很快就会因为只能“求和”取实值的函数而感到不便；在许多情形，懂得如何去定义定义在 X 上、取值于 R 上的拓扑向量空间（例如一个 Banach 空间或 Banach 空间上的算子空间）中的函数的积分，是有用的。人们可以确认，这一延拓容易实现，无需对积分理论作任何深刻的修改。

在上述概述中，我们让连续函数扮演了占优势的角色；自然要问：测度的概念是否真的本质地依赖于它在其上定义的集合 X 上一个拓扑的存在。对理论的细致考察表明完全不是这么回事，而且延拓方法同样适用于定义在向量空间 V 上的正线性形式 \(\mu(f)\) ，这里 V 由定义在任意集合 X 上的数值函数组成，只需对 V 和 \(\mu(f)\) 附加某些补充条件；当 V 是具有紧支集的连续函数所成的空间 \(\mathcal{K}(\mathrm{X})\) 时，这些条件自动满足，但它们在更一般的情形中也能满足。不过，这种更大的普遍性在某些方面是虚幻的：事实上可以证明，每个“抽象测度”在某种意义下都“同构”于（从连续函数出发）定义在一个适当的局部紧空间上的测度；另一方面，大多数应用所涉及的集合 X 都带有在该问题中自然出现的拓扑；因此，直到第 IX 章，我们将只研究定义在局部紧空间上的测度。

前两章是本理论的预备篇章：它们致力于证明一些对延拓具有根本意义的不等式，并研究某些序向量空间，即 Riesz 空间，这些空间在后面的若干问题中起着重要的作用。

局部紧空间上的测度概念在第 III 章中定义；我们以 Riesz 定理为出发点，于是该定理在这里成为一个定义：连续函数的积分先于集合的测度得到定义，即定义为 \(\mathcal{K}(\mathrm{X})\) 上的一个正线性形式。这种表述方式具有某些技术上的优点（特别是因为连续函数构成一个向量空间，而集合的特征函数却不然）；此外，正是在 \(\mathcal{K}(\mathrm{X})\) 上的泛函的形式下，积分在大量问题中自然地出现。最后，\(\mathcal{K}(\mathrm{X})\) 上两个正线性形式之差（我们仍称之为 X 上的测度）可以刻画为 \(\mathcal{K}(\mathrm{X})\) 上满足某些连续性条件的线性形式；这样，积分理论一方面联系于拓扑向量空间中的一般对偶理论（参见第 V 卷），另一方面联系于分布理论，分布推广了测度概念的某些方面，我们将在随后的一卷中陈述它。

第 IV 章专门讨论积分的延拓；在那里同时定义了可积函数与集合的测度，还定义了函数空间 \(L^{p}\) ，它们在应用中的重要性是可观的；那里还说明了，可测函数概念的引入如何导出便于使用的可积性判别法。

在其后两章中将看到，可测函数又如何作为“密度”出现，使得能从一个给定的测度出发，在空间 X 上定义新的测度。这项研究特别导出 \(L^{p}\) 空间对偶理论中的若干重要结果，它还与向量测度的概念相联系；在最有利的情形，向量测度可以归结为向量值函数（关于一个正测度）的积分理论。

我们还将展开可以视为积分学创立者们所引入的平面面积与体积“切片分解”思想的现代极致的内容：在一定条件下，空间 X 上的一个测度可以分解为一个“和”，和中的每个测度都由空间 X 的一条“切片”（也就是关于某个关系 R 的一个等价类）所承载；此外，这样的分解使得能够通过先“在每条切片上”积分、再（关于一个适当的测度）对商空间 X/R 上由此得到的函数积分，来计算一个函数关于原来那个测度的积分（这是对指标取遍一个乘积集的族的和式的“二重求和”的一个推广）。

第 VII 章专门研究局部紧群上的 Haar 测度，它在相差一个标量倍数的意义下，由在群的每个左平移下不变这一性质所刻画。

第 VIII 章陈述测度的卷积，这一概念在现代泛函分析中起着头等重要的作用。

第 IX 章专门讨论未必局部紧的 Hausdorff 拓扑空间中的积分，特别是局部凸向量空间中的积分；这特别使得 Fourier 变换的理论能够延拓到后面这类空间。开头几节所选的叙述方式，在于尽可能化归为前面各章中已处理过的紧空间的情形。

